\begin{abstract}
We formalize the pathwise correspondence between the Reed--Frost (Independent Cascade) epidemic
and bond percolation, after Kempe, Kleinberg and Tardos: with one coin per edge, the nodes
infected in round $t$ are exactly those at distance $t$ from the initial set in the graph of
open edges. Consequently the final outbreak is the set of nodes connected to the initial set by
open edges, and with i.i.d.\ Bernoulli$(p)$ coins the probability of eventual infection is a
bond-percolation connection probability. We also formalize the duality between the
coalescing-branching random walk (COBRA) and the BIPS epidemic of Cooper, Radzik and Rivera
(Theorem 4), pathwise and by time reversal of i.i.d.\ rounds. The development is in the finite
layer of the shared \texttt{dynamics} package.
\end{abstract}

\section{The model}
\begin{definition}[Percolated graph and epidemic]\label{def:reedfrost}
\lean{Epidemics.perc,Epidemics.SIR,Epidemics.step,Epidemics.run,Epidemics.setDist}\leanok
Every edge $e$ of a finite graph $G$ carries a coin $\omega(e)\in\{\mathrm{true},\mathrm{false}\}$;
the percolated graph $G_\omega$ keeps the edges whose coin is true. In one round, a susceptible node
becomes infected if it has an infected neighbour across an open edge, and every infected node
recovers. Initially the set $I_0$ is infected and nobody is recovered. We write $d(v)\in\mathbb N\cup\{\infty\}$
for the distance from $I_0$ to $v$ in $G_\omega$.
\end{definition}

\section{Distances from a set}
\begin{lemma}[Distance facts]\label{lem:dist}
\lean{Epidemics.setDist_eq_subtype,Epidemics.setDist_eq_zero_iff,Epidemics.setDist_le_add_one,Epidemics.exists_edist_eq_setDist,Epidemics.exists_adj_of_setDist_succ,Epidemics.setDist_ne_top_iff,Epidemics.setDist_lt_card}\leanok\uses{def:reedfrost}
$d(v)=0$ iff $v\in I_0$; if $u$ and $v$ are adjacent in $G_\omega$ then $d(v)\le d(u)+1$; if
$d(v)=t+1$ then $v$ has a neighbour $u$ in $G_\omega$ with $d(u)=t$; $d(v)<\infty$ iff $v$ is
connected to $I_0$ in $G_\omega$, and then $d(v)<|V|$.
\end{lemma}
\begin{proof}\leanok
The infimum over $I_0$ is attained; take the second-to-last vertex of a shortest walk; a shortest
walk is a path.
\end{proof}

\section{Pathwise layers}
\begin{theorem}[Infected and recovered sets]\label{thm:layers}
\lean{Epidemics.infected_recovered_iff,Epidemics.infected_iff,Epidemics.recovered_iff}\leanok\uses{lem:dist}
For every coin assignment and every $t$, the nodes infected in round $t$ are those with $d(v)=t$,
and the nodes recovered by round $t$ are those with $d(v)<t$.
\end{theorem}
\begin{proof}\leanok
Both statements together, by induction on $t$, using Lemma~\ref{lem:dist}.
\end{proof}
\begin{corollary}[Extinction and final size]\label{cor:final}
\lean{Epidemics.no_infected_of_dist_lt,Epidemics.extinct_of_dist_lt,Epidemics.extinct,Epidemics.final_recovered_iff}\leanok\uses{thm:layers}
Nobody is infected in round $t$ as soon as every finite $d(v)$ is $<t$, in particular for
$t\ge|V|$; the nodes recovered after $|V|$ rounds are those connected to $I_0$ in $G_\omega$.
\end{corollary}
\begin{proof}\leanok
Immediate from Theorem~\ref{thm:layers} and $d(v)<|V|$ when finite.
\end{proof}

\section{Random coins}
\begin{theorem}[Reed--Frost $\Leftrightarrow$ bond percolation]\label{thm:perc}
\lean{Epidemics.bernoulli,Epidemics.coins,Epidemics.coins_prob_open,Epidemics.prob_infected_eq_prob_connected}\leanok\uses{cor:final}
With independent Bernoulli$(p)$ coins, each edge is open with probability $p$, and the
probability that $v$ is eventually infected equals the probability that $v$ is connected to $I_0$
by open edges.
\end{theorem}
\begin{proof}\leanok
The two events coincide for every coin assignment (Corollary~\ref{cor:final}).
\end{proof}
The coupling is pathwise because every edge is used at most once: once its coin has been seen
from an infected endpoint, that endpoint is recovered for good. This is why a single coin per edge
gives the same process, in distribution, as fresh coins in every round (that distributional
statement is not formalized here).

\section{COBRA and BIPS: the model}
We follow Cooper, Radzik and Rivera, \emph{The coalescing-branching random walk on expanders and
the dual epidemic process}, PODC 2016 (arXiv:1602.05768).
\begin{definition}[Rounds of neighbour choices]\label{def:choices}
\lean{Epidemics.Choices,Epidemics.instFintypeChoices,Epidemics.choices_nonempty}\leanok
In a round $r$, every vertex $x$ of a finite graph $G$ samples $k$ neighbours
$r(x,0),\dots,r(x,k-1)$. A uniform round is exactly the paper's sampling: $k$ uniform neighbours
with replacement, independently for every vertex; $t$ i.i.d.\ uniform rounds are averaged with
\texttt{Dynamics.expList}. If $G$ is connected and has at least two vertices, rounds exist.
\end{definition}
\begin{proof}\leanok
Every vertex of a connected graph with at least two vertices has a neighbour.
\end{proof}
\begin{definition}[COBRA and BIPS]\label{def:cobra}
\lean{Epidemics.cobraStep,Epidemics.cobraRun,Epidemics.bipsStep,Epidemics.bipsRun}\leanok\uses{def:choices}
COBRA from $C_0=C$: $C_{s+1}=\{r_{s+1}(x,i) : x\in C_s,\ i<k\}$. BIPS with persistent source $v$
from $A_0$: $A_{s+1}=\{v\}\cup\{u : \exists i,\ r_{s+1}(u,i)\in A_s\}$.
\end{definition}

\section{COBRA--BIPS duality}
\begin{lemma}[Time reversal of i.i.d.\ rounds]\label{lem:reverse}
\lean{Epidemics.expList_reverse,Epidemics.expList_congr_length}\leanok
For every functional $F$ of $T$ i.i.d.\ uniform rounds,
$\mathbb E[F(r_T,\dots,r_1)]=\mathbb E[F(r_1,\dots,r_T)]$.
\end{lemma}
\begin{proof}\leanok
The uniform average over $\mathrm{Fin}\,T\to\alpha$ is invariant under $\omega\mapsto\omega\circ\mathrm{rev}$.
\end{proof}
\begin{lemma}[Pathwise duality]\label{lem:pathwise}
\lean{Epidemics.cobra_hit_iff_bips_reverse}\leanok\uses{def:cobra}
For all rounds $r_1,\dots,r_t$: COBRA from $C$ along $r_1,\dots,r_t$ visits $v$ at some time
$s\le t$ (time $0$ included) iff BIPS from $\{v\}$ along $r_t,\dots,r_1$ infects a vertex of $C$.
\end{lemma}
\begin{proof}\leanok
Induction on $t$, generalizing $C$. Both sides say that a chain $x_0\in C$,
$x_j=r_j(x_{j-1},i_j)$, reaches $v$ within $t$ steps: COBRA unfolds the first round, BIPS the
last round of the reversed sequence.
\end{proof}
\begin{theorem}[Theorem 4 of Cooper--Radzik--Rivera]\label{thm:duality}
\lean{Epidemics.cobra_bips_duality,Epidemics.cobra_bips_duality_singleton}\leanok\uses{lem:reverse,lem:pathwise}
For every finite graph, every $k$, every vertex $v$, set $C$ and time $t$:
$\hat{\mathbb P}(\mathrm{Hit}_C(v)>t\mid C_0=C)=\mathbb P(C\cap A_t=\emptyset\mid A_0=\{v\})$. In
particular (equation (2)), $\hat{\mathbb P}(\mathrm{Hit}_u(v)>t)=\mathbb P(u\notin A_t\mid A_0=\{v\})$.
The paper assumes $G$ connected and regular and $k\ge1$; none of this is needed.
\end{theorem}
\begin{proof}\leanok
Lemma~\ref{lem:pathwise} identifies the two events after reversing the rounds, and
Lemma~\ref{lem:reverse} says that reversing the rounds preserves their law. The paper proves the
theorem instead by induction on $t$ with one-step conditioning.
\end{proof}
